
\documentclass{article}
\usepackage{colortbl}
\usepackage{makecell}
\usepackage{multirow}
\usepackage{supertabular}

\begin{document}

\newcounter{utterance}

\centering \large Interaction Transcript for game `clean\_up', experiment `2\_hard\_3obj\_ro', episode 3 with gemma{-}4{-}26b{-}a4b{-}it.
\vspace{24pt}

{ \footnotesize  \setcounter{utterance}{1}
\setlength{\tabcolsep}{0pt}
\begin{supertabular}{c@{$\;$}|p{.15\linewidth}@{}p{.15\linewidth}p{.15\linewidth}p{.15\linewidth}p{.15\linewidth}p{.15\linewidth}}
    \# & $\;$A & \multicolumn{4}{c}{Game Master} & $\;\:$B\\
    \hline

    \theutterance \stepcounter{utterance}  
    & & \multicolumn{4}{p{0.6\linewidth}}{
        \cellcolor[rgb]{0.9,0.9,0.9}{
            \makecell[{{p{\linewidth}}}]{
                \texttt{\tiny{[P1$\langle$GM]}}
                \texttt{Eu sunt maestrul vostru de joc, iar voi jucați un joc colaborativ cu următoarea grilă ca tablă de joc:} \\
\\ 
\texttt{```} \\
\texttt{ 1234567} \\
\texttt{╔═╤═╤═══╗} \\
\texttt{║L│◌│◌◌◌║ 1} \\
\texttt{║C├─┤◌◌◌║ 2} \\
\texttt{╟─┤◌├─┬─╢ 3} \\
\texttt{║◌│◌│◌│◌║ 4} \\
\texttt{╟─┼─┘◌│◌║ 5} \\
\texttt{║◌│◌◌┌┴─╢ 6} \\
\texttt{║◌│◌◌│◌P║ 7} \\
\texttt{╚═╧══╧══╝} \\
\\ 
\texttt{```} \\
\\ 
\texttt{* Marginea superioară afișează coordonatele x crescând spre dreapta, iar marginea dreaptă coordonatele y crescând în jos.} \\
\texttt{* Următoarele obiecte sunt plasate aleatoriu pe grila voastră: 'C', 'L', 'P'.} \\
\\ 
\\ 
\texttt{Celălalt jucător vede o variantă a tablei de joc, în care obiectele sunt plasate în diferite locații aleatorii. Nu puteți vedea tabla de joc a celuilalt jucător, iar el nu o poate vedea pe a dumneavoastră.} \\
\\ 
\\ 
\texttt{**Obiectiv:**} \\
\\ 
\texttt{Ambii jucători trebuie să mute obiectele de pe fundalul lor respectiv astfel încât obiectele identice să ajungă la aceleași coordonate. Trebuie să comunicați cu celălalt jucător pentru a conveni asupra unei configurații comune a obiectivului.} \\
\\ 
\\ 
\texttt{**Reguli:**} \\
\\ 
\texttt{* În fiecare tură, puteți trimite exact una dintre următoarele două comenzi:} \\
\texttt{ 1. `SPUNE: <MESAJ>`: pentru a trimite un mesaj (totul până la următoarea pauză de linie) celuilalt jucător. Eu îl voi transmite partenerului tău.} \\
\texttt{2. `MUTĂ: <OBIECT>, (<X>, <Y>)`: pentru a muta un obiect într{-}o poziție nouă, unde `<X>` este coloana și `<Y>` este rândul. Te voi informa dacă mutarea ta a avut succes sau nu.} \\
\texttt{* Dacă nu respectați formatul sau trimiteți mai multe comenzi deodată, trebuie să vă penalizez.} \\
\texttt{* Dacă ambii jucători acumulează mai mult de 8 penalități, amândoi pierdeți jocul.} \\
\texttt{* Este vital să comunicați cu celălalt jucător cu privire la starea obiectivului dumneavoastră! *Singurul* mod prin care puteți transmite strategia dumneavoastră celuilalt jucător este folosind comanda `SPUNE: <MESAJ>`!} \\
\\ 
\\ 
\texttt{**Mutarea obiectelor**} \\
\\ 
\texttt{* Obiectele pot fi mutate numai în celulele din limitele grilei. Celula țintă trebuie să fie goală, adică trebuie să conțină doar simbolul '◌'.} \\
\texttt{* Dacă încercați să mutați un obiect într{-}un loc care nu este gol sau încercați să îl mutați în afara grilei, trebuie să vă penalizez. Aveți dreptul la o nouă încercare.} \\
\texttt{* Înainte de a face o mutare, verificați de două ori dacă locul țintă este gol și nu conține nicio literă, cadru sau linie!} \\
\\ 
\\ 
\texttt{**Sfârșitul jocului**} \\
\\ 
\texttt{Dacă credeți că ați atins obiectivul de aliniere a tuturor obiectelor, puteți cere celuilalt jucător să încheie jocul trimițând `SPUNE: terminat?`. Dacă celălalt jucător vă cere să încheiați jocul, iar dvs. răspundeți cu `SPUNE: terminat!`, jocul se va încheia.} \\
\\ 
\texttt{Ambii jucători câștigă dacă jocul se încheie în decursul a 12 runde, unde o rundă este definită ca fiind trimiterea de către doi jucători a câte unei comenzi valide.} \\
\\ 
\texttt{**Punctaj:**} \\
\\ 
\texttt{Cu cât obiectele identice sunt mai apropiate pe ambele table de joc, cu atât primești mai multe puncte. Penalizările îți reduc punctele. Poți bate recordul?} \\
\\ 
\\ 
\texttt{Vă rugăm să trimiteți un mesaj celuilalt jucător pentru a începe jocul!} \\
            }
        }
    }
    & \\ \\

    \theutterance \stepcounter{utterance}  
    & \multicolumn{4}{p{0.6\linewidth}}{
        \cellcolor[rgb]{0.9,0.9,0.9}{
            \makecell[{{p{\linewidth}}}]{
                \texttt{\tiny{[P1$\rangle$GM]}}
                \texttt{SPUNE: Salut! Am început jocul. Pe tabla mea, obiectele sunt la următoarele coordonate: L(1,1), C(1,2) și P(7,7). Hai să comunicăm locațiile noastre pentru a găsi o configurație comună. Spune{-}mi unde sunt obiectele tale!} \\
            }
        }
    }
    & & \\ \\

    \theutterance \stepcounter{utterance}  
    & & & \multicolumn{2}{p{0.3\linewidth}}{
        \cellcolor[rgb]{0.9,0.9,0.9}{
            \makecell[{{p{\linewidth}}}]{
                \texttt{\tiny{[GM$|$GM]}}
                \texttt{SPUNE: Salut! Am început jocul. Pe tabla mea, obiectele sunt la următoarele coordonate: L(1,1), C(1,2) și P(7,7). Hai să comunicăm locațiile noastre pentru a găsi o configurație comună. Spune{-}mi unde sunt obiectele tale!} \\
            }
        }
    }
    & & \\ \\

    \theutterance \stepcounter{utterance}  
    & & \multicolumn{4}{p{0.6\linewidth}}{
        \cellcolor[rgb]{0.9,0.9,0.9}{
            \makecell[{{p{\linewidth}}}]{
                \texttt{\tiny{[GM$\rangle$P2]}}
                \texttt{Eu sunt maestrul vostru de joc, iar voi jucați un joc colaborativ cu următoarea grilă ca tablă de joc:} \\
\\ 
\texttt{```} \\
\texttt{ 1234567} \\
\texttt{╔═╤═╤═══╗} \\
\texttt{║◌│◌│◌◌◌║ 1} \\
\texttt{║◌├─┤◌L◌║ 2} \\
\texttt{╟─┤P├─┬─╢ 3} \\
\texttt{║◌│◌│◌│◌║ 4} \\
\texttt{╟─┼─┘C│◌║ 5} \\
\texttt{║◌│◌◌┌┴─╢ 6} \\
\texttt{║◌│◌◌│◌◌║ 7} \\
\texttt{╚═╧══╧══╝} \\
\\ 
\texttt{```} \\
\\ 
\texttt{* Marginea superioară afișează coordonatele x crescând spre dreapta, iar marginea dreaptă coordonatele y crescând în jos.} \\
\texttt{* Următoarele obiecte sunt plasate aleatoriu pe grila voastră: 'C', 'L', 'P'.} \\
\\ 
\\ 
\texttt{Celălalt jucător vede o variantă a tablei de joc, în care obiectele sunt plasate în diferite locații aleatorii. Nu puteți vedea tabla de joc a celuilalt jucător, iar el nu o poate vedea pe a dumneavoastră.} \\
\\ 
\\ 
\texttt{**Obiectiv:**} \\
\\ 
\texttt{Ambii jucători trebuie să mute obiectele de pe fundalul lor respectiv astfel încât obiectele identice să ajungă la aceleași coordonate. Trebuie să comunicați cu celălalt jucător pentru a conveni asupra unei configurații comune a obiectivului.} \\
\\ 
\\ 
\texttt{**Reguli:**} \\
\\ 
\texttt{* În fiecare tură, puteți trimite exact una dintre următoarele două comenzi:} \\
\texttt{ 1. `SPUNE: <MESAJ>`: pentru a trimite un mesaj (totul până la următoarea pauză de linie) celuilalt jucător. Eu îl voi transmite partenerului tău.} \\
\texttt{2. `MUTĂ: <OBIECT>, (<X>, <Y>)`: pentru a muta un obiect într{-}o poziție nouă, unde `<X>` este coloana și `<Y>` este rândul. Te voi informa dacă mutarea ta a avut succes sau nu.} \\
\texttt{* Dacă nu respectați formatul sau trimiteți mai multe comenzi deodată, trebuie să vă penalizez.} \\
\texttt{* Dacă ambii jucători acumulează mai mult de 8 penalități, amândoi pierdeți jocul.} \\
\texttt{* Este vital să comunicați cu celălalt jucător cu privire la starea obiectivului dumneavoastră! *Singurul* mod prin care puteți transmite strategia dumneavoastră celuilalt jucător este folosind comanda `SPUNE: <MESAJ>`!} \\
\\ 
\\ 
\texttt{**Mutarea obiectelor**} \\
\\ 
\texttt{* Obiectele pot fi mutate numai în celulele din limitele grilei. Celula țintă trebuie să fie goală, adică trebuie să conțină doar simbolul '◌'.} \\
\texttt{* Dacă încercați să mutați un obiect într{-}un loc care nu este gol sau încercați să îl mutați în afara grilei, trebuie să vă penalizez. Aveți dreptul la o nouă încercare.} \\
\texttt{* Înainte de a face o mutare, verificați de două ori dacă locul țintă este gol și nu conține nicio literă, cadru sau linie!} \\
\\ 
\\ 
\texttt{**Sfârșitul jocului**} \\
\\ 
\texttt{Dacă credeți că ați atins obiectivul de aliniere a tuturor obiectelor, puteți cere celuilalt jucător să încheie jocul trimițând `SPUNE: terminat?`. Dacă celălalt jucător vă cere să încheiați jocul, iar dvs. răspundeți cu `SPUNE: terminat!`, jocul se va încheia.} \\
\\ 
\texttt{Ambii jucători câștigă dacă jocul se încheie în decursul a 12 runde, unde o rundă este definită ca fiind trimiterea de către doi jucători a câte unei comenzi valide.} \\
\\ 
\texttt{**Punctaj:**} \\
\\ 
\texttt{Cu cât obiectele identice sunt mai apropiate pe ambele table de joc, cu atât primești mai multe puncte. Penalizările îți reduc punctele. Poți bate recordul?} \\
\\ 
\texttt{Celălalt jucător a început jocul trimițând acest mesaj:} \\
\\ 
\texttt{"Salut! Am început jocul. Pe tabla mea, obiectele sunt la următoarele coordonate: L(1,1), C(1,2) și P(7,7). Hai să comunicăm locațiile noastre pentru a găsi o configurație comună. Spune{-}mi unde sunt obiectele tale!"} \\
\\ 
\texttt{Care este prima ta comandă?} \\
            }
        }
    }
    & \\ \\

    \theutterance \stepcounter{utterance}  
    & & & \multicolumn{4}{p{0.6\linewidth}}{
        \cellcolor[rgb]{0.9,0.9,0.9}{
            \makecell[{{p{\linewidth}}}]{
                \texttt{\tiny{[GM$\langle$P2]}}
                \texttt{SPUNE: Salut! Am înțeles. Pe tabla mea, obiectele se află la următoarele coordonate: L(5,2), P(3,3) și C(5,5). Trebuie să găsim o locație comună pentru toate. Ce coordonate propui pentru L, C și P?} \\
            }
        }
    }
     \\ \\

    \theutterance \stepcounter{utterance}  
    & & & \multicolumn{2}{p{0.3\linewidth}}{
        \cellcolor[rgb]{0.9,0.9,0.9}{
            \makecell[{{p{\linewidth}}}]{
                \texttt{\tiny{[GM$|$GM]}}
                \texttt{SPUNE: Salut! Am înțeles. Pe tabla mea, obiectele se află la următoarele coordonate: L(5,2), P(3,3) și C(5,5). Trebuie să găsim o locație comună pentru toate. Ce coordonate propui pentru L, C și P?} \\
            }
        }
    }
    & & \\ \\

    \theutterance \stepcounter{utterance}  
    & & \multicolumn{4}{p{0.6\linewidth}}{
        \cellcolor[rgb]{0.9,0.9,0.9}{
            \makecell[{{p{\linewidth}}}]{
                \texttt{\tiny{[P1$\langle$GM]}}
                \texttt{Mesajul dvs. a fost transmis celuilalt jucător.} \\
\\ 
\\ 
\texttt{În prezent jucați runda 1 din maximum 12.} \\
\texttt{Ați acumulat în mod colectiv 0 din 8 penalizări.} \\
\texttt{Celălalt jucător a trimis acest mesaj:} \\
\\ 
\\ 
\texttt{"Salut! Am înțeles. Pe tabla mea, obiectele se află la următoarele coordonate: L(5,2), P(3,3) și C(5,5). Trebuie să găsim o locație comună pentru toate. Ce coordonate propui pentru L, C și P?"} \\
\\ 
\\ 
\texttt{Care este următoarea dvs. comandă?} \\
            }
        }
    }
    & \\ \\

    \theutterance \stepcounter{utterance}  
    & \multicolumn{4}{p{0.6\linewidth}}{
        \cellcolor[rgb]{0.9,0.9,0.9}{
            \makecell[{{p{\linewidth}}}]{
                \texttt{\tiny{[P1$\rangle$GM]}}
                \texttt{SPUNE: Salut! Mulțumesc pentru informații. Avem:} \\
\texttt{Eu: L(1,1), C(1,2), P(7,7)} \\
\texttt{Tu: L(5,2), P(3,3), C(5,5)} \\
\\ 
\texttt{Ce zici să încercăm să le aducem pe toate într{-}un colț sau pe o linie? De exemplu, putem încerca să le punem pe L(2,2), C(2,3) și P(2,4)? Sau ai o altă preferință care să fie ușor de realizat pentru amândoi?} \\
            }
        }
    }
    & & \\ \\

    \theutterance \stepcounter{utterance}  
    & & & \multicolumn{2}{p{0.3\linewidth}}{
        \cellcolor[rgb]{0.9,0.9,0.9}{
            \makecell[{{p{\linewidth}}}]{
                \texttt{\tiny{[GM$|$GM]}}
                \texttt{SPUNE: Salut! Mulțumesc pentru informații. Avem:} \\
\texttt{Eu: L(1,1), C(1,2), P(7,7)} \\
\texttt{Tu: L(5,2), P(3,3), C(5,5)} \\
\\ 
\texttt{Ce zici să încercăm să le aducem pe toate într{-}un colț sau pe o linie? De exemplu, putem încerca să le punem pe L(2,2), C(2,3) și P(2,4)? Sau ai o altă preferință care să fie ușor de realizat pentru amândoi?} \\
            }
        }
    }
    & & \\ \\

    \theutterance \stepcounter{utterance}  
    & & & \multicolumn{2}{p{0.3\linewidth}}{
        \cellcolor[rgb]{0.9,0.9,0.9}{
            \makecell[{{p{\linewidth}}}]{
                \texttt{\tiny{[GM$|$GM]}}
                \texttt{Invalid format: tail is not empty: '} \\
\texttt{Eu: L(1,1), C(1,2), P(7,7)} \\
\texttt{Tu: L(5,2), P(3,3), C(5,5)} \\
\\ 
\texttt{Ce zici să încercăm să le aducem pe toate într{-}un colț sau pe o linie? De exemplu, putem încerca să le punem pe L(2,2), C(2,3) și P(2,4)? Sau ai o altă preferință care să fie ușor de realizat pentru amândoi?} \\
            }
        }
    }
    & & \\ \\

    \theutterance \stepcounter{utterance}  
    & & \multicolumn{4}{p{0.6\linewidth}}{
        \cellcolor[rgb]{0.9,0.9,0.9}{
            \makecell[{{p{\linewidth}}}]{
                \texttt{\tiny{[P1$\langle$GM]}}
                \texttt{Penalizare: Mesajul dvs. nu trebuie să conțină nimic după comandă!} \\
\\ 
\texttt{Asigurați{-}vă că răspunsul dvs. conține doar `SPUNE: <MESAJ>` sau `MUTĂ: <OBIECT>, (<X>, <Y>)`, și nimic altceva!} \\
\\ 
\texttt{Ați acumulat în mod colectiv 1 din 8 penalizări.} \\
\texttt{Vă rugăm să încercați din nou!} \\
            }
        }
    }
    & \\ \\

    \theutterance \stepcounter{utterance}  
    & \multicolumn{4}{p{0.6\linewidth}}{
        \cellcolor[rgb]{0.9,0.9,0.9}{
            \makecell[{{p{\linewidth}}}]{
                \texttt{\tiny{[P1$\rangle$GM]}}
                \texttt{SPUNE: Ce zici să încercăm să le aducem pe toate pe coloana 2? Propun L(2,2), C(2,3) și P(2,4). Este o locație bună pentru amândoi?} \\
            }
        }
    }
    & & \\ \\

    \theutterance \stepcounter{utterance}  
    & & & \multicolumn{2}{p{0.3\linewidth}}{
        \cellcolor[rgb]{0.9,0.9,0.9}{
            \makecell[{{p{\linewidth}}}]{
                \texttt{\tiny{[GM$|$GM]}}
                \texttt{SPUNE: Ce zici să încercăm să le aducem pe toate pe coloana 2? Propun L(2,2), C(2,3) și P(2,4). Este o locație bună pentru amândoi?} \\
            }
        }
    }
    & & \\ \\

    \theutterance \stepcounter{utterance}  
    & & \multicolumn{4}{p{0.6\linewidth}}{
        \cellcolor[rgb]{0.9,0.9,0.9}{
            \makecell[{{p{\linewidth}}}]{
                \texttt{\tiny{[GM$\rangle$P2]}}
                \texttt{Mesajul dvs. a fost transmis celuilalt jucător.} \\
\\ 
\\ 
\texttt{În prezent jucați runda 2 din maximum 12.} \\
\texttt{Ați acumulat în mod colectiv 1 din 8 penalizări.} \\
\texttt{Celălalt jucător a trimis acest mesaj:} \\
\\ 
\\ 
\texttt{"Ce zici să încercăm să le aducem pe toate pe coloana 2? Propun L(2,2), C(2,3) și P(2,4). Este o locație bună pentru amândoi?"} \\
\\ 
\\ 
\texttt{Care este următoarea dvs. comandă?} \\
            }
        }
    }
    & \\ \\

    \theutterance \stepcounter{utterance}  
    & & & \multicolumn{4}{p{0.6\linewidth}}{
        \cellcolor[rgb]{0.9,0.9,0.9}{
            \makecell[{{p{\linewidth}}}]{
                \texttt{\tiny{[GM$\langle$P2]}}
                \texttt{SPUNE: Propunerea ta sună bine! Să verificăm dacă locațiile sunt libere pe tabla mea. L(2,2) este liber, C(2,3) este liber și P(2,4) este liber. Pe tabla ta, L(2,2), C(2,3) și P(2,4) sunt libere? Dacă da, putem începe mutările.} \\
            }
        }
    }
     \\ \\

    \theutterance \stepcounter{utterance}  
    & & & \multicolumn{2}{p{0.3\linewidth}}{
        \cellcolor[rgb]{0.9,0.9,0.9}{
            \makecell[{{p{\linewidth}}}]{
                \texttt{\tiny{[GM$|$GM]}}
                \texttt{SPUNE: Propunerea ta sună bine! Să verificăm dacă locațiile sunt libere pe tabla mea. L(2,2) este liber, C(2,3) este liber și P(2,4) este liber. Pe tabla ta, L(2,2), C(2,3) și P(2,4) sunt libere? Dacă da, putem începe mutările.} \\
            }
        }
    }
    & & \\ \\

    \theutterance \stepcounter{utterance}  
    & & \multicolumn{4}{p{0.6\linewidth}}{
        \cellcolor[rgb]{0.9,0.9,0.9}{
            \makecell[{{p{\linewidth}}}]{
                \texttt{\tiny{[P1$\langle$GM]}}
                \texttt{Mesajul dvs. a fost transmis celuilalt jucător.} \\
\\ 
\\ 
\texttt{În prezent jucați runda 2 din maximum 12.} \\
\texttt{Ați acumulat în mod colectiv 1 din 8 penalizări.} \\
\texttt{Celălalt jucător a trimis acest mesaj:} \\
\\ 
\\ 
\texttt{"Propunerea ta sună bine! Să verificăm dacă locațiile sunt libere pe tabla mea. L(2,2) este liber, C(2,3) este liber și P(2,4) este liber. Pe tabla ta, L(2,2), C(2,3) și P(2,4) sunt libere? Dacă da, putem începe mutările."} \\
\\ 
\\ 
\texttt{Care este următoarea dvs. comandă?} \\
            }
        }
    }
    & \\ \\

    \theutterance \stepcounter{utterance}  
    & \multicolumn{4}{p{0.6\linewidth}}{
        \cellcolor[rgb]{0.9,0.9,0.9}{
            \makecell[{{p{\linewidth}}}]{
                \texttt{\tiny{[P1$\rangle$GM]}}
                \texttt{SPUNE: Pe tabla mea, L(2,2), C(2,3) și P(2,4) sunt libere. Să începem! Eu mut prima: MUTĂ: L, (2, 2)} \\
            }
        }
    }
    & & \\ \\

    \theutterance \stepcounter{utterance}  
    & & & \multicolumn{2}{p{0.3\linewidth}}{
        \cellcolor[rgb]{0.9,0.9,0.9}{
            \makecell[{{p{\linewidth}}}]{
                \texttt{\tiny{[GM$|$GM]}}
                \texttt{SPUNE: Pe tabla mea, L(2,2), C(2,3) și P(2,4) sunt libere. Să începem! Eu mut prima: MUTĂ: L, (2, 2)} \\
            }
        }
    }
    & & \\ \\

    \theutterance \stepcounter{utterance}  
    & & & \multicolumn{2}{p{0.3\linewidth}}{
        \cellcolor[rgb]{0.9,0.9,0.9}{
            \makecell[{{p{\linewidth}}}]{
                \texttt{\tiny{[GM$|$GM]}}
                \texttt{Invalid response format: SPUNE: Pe tabla mea, L(2,2), C(2,3) și P(2,4) sunt libere. Să începem! Eu mut prima: MUTĂ: L, (2, 2)} \\
            }
        }
    }
    & & \\ \\

    \theutterance \stepcounter{utterance}  
    & & \multicolumn{4}{p{0.6\linewidth}}{
        \cellcolor[rgb]{0.9,0.9,0.9}{
            \makecell[{{p{\linewidth}}}]{
                \texttt{\tiny{[P1$\langle$GM]}}
                \texttt{Penalizare: Mesajul dvs. conține mai mult de o comandă!} \\
\\ 
\texttt{Asigurați{-}vă că răspunsul dvs. conține doar `SPUNE: <MESAJ>` sau `MUTĂ: <OBIECT>, (<X>, <Y>)`, și nimic altceva!} \\
\\ 
\texttt{Ați acumulat în mod colectiv 2 din 8 penalizări.} \\
\texttt{Vă rugăm să încercați din nou!} \\
            }
        }
    }
    & \\ \\

    \theutterance \stepcounter{utterance}  
    & \multicolumn{4}{p{0.6\linewidth}}{
        \cellcolor[rgb]{0.9,0.9,0.9}{
            \makecell[{{p{\linewidth}}}]{
                \texttt{\tiny{[P1$\rangle$GM]}}
                \texttt{MUTĂ: L, (2, 2)} \\
            }
        }
    }
    & & \\ \\

    \theutterance \stepcounter{utterance}  
    & & & \multicolumn{2}{p{0.3\linewidth}}{
        \cellcolor[rgb]{0.9,0.9,0.9}{
            \makecell[{{p{\linewidth}}}]{
                \texttt{\tiny{[GM$|$GM]}}
                \texttt{MUTĂ: L, (2, 2)} \\
            }
        }
    }
    & & \\ \\

    \theutterance \stepcounter{utterance}  
    & & & \multicolumn{2}{p{0.3\linewidth}}{
        \cellcolor[rgb]{0.9,0.9,0.9}{
            \makecell[{{p{\linewidth}}}]{
                \texttt{\tiny{[GM$|$GM]}}
                \texttt{Penalizare: (2,2) nu este gol, ci conține '├'.} \\
\\ 
\texttt{Ați acumulat în mod colectiv 3 din 8 penalizări.} \\
\texttt{Vă rugăm să încercați din nou!} \\
            }
        }
    }
    & & \\ \\

    \theutterance \stepcounter{utterance}  
    & & \multicolumn{4}{p{0.6\linewidth}}{
        \cellcolor[rgb]{0.9,0.9,0.9}{
            \makecell[{{p{\linewidth}}}]{
                \texttt{\tiny{[P1$\langle$GM]}}
                \texttt{Penalizare: (2,2) nu este gol, ci conține '├'.} \\
\\ 
\texttt{Ați acumulat în mod colectiv 3 din 8 penalizări.} \\
\texttt{Vă rugăm să încercați din nou!} \\
            }
        }
    }
    & \\ \\

    \theutterance \stepcounter{utterance}  
    & \multicolumn{4}{p{0.6\linewidth}}{
        \cellcolor[rgb]{0.9,0.9,0.9}{
            \makecell[{{p{\linewidth}}}]{
                \texttt{\tiny{[P1$\rangle$GM]}}
                \texttt{MUTĂ: L, (2, 1)} \\
            }
        }
    }
    & & \\ \\

    \theutterance \stepcounter{utterance}  
    & & & \multicolumn{2}{p{0.3\linewidth}}{
        \cellcolor[rgb]{0.9,0.9,0.9}{
            \makecell[{{p{\linewidth}}}]{
                \texttt{\tiny{[GM$|$GM]}}
                \texttt{MUTĂ: L, (2, 1)} \\
            }
        }
    }
    & & \\ \\

    \theutterance \stepcounter{utterance}  
    & & & \multicolumn{2}{p{0.3\linewidth}}{
        \cellcolor[rgb]{0.9,0.9,0.9}{
            \makecell[{{p{\linewidth}}}]{
                \texttt{\tiny{[GM$|$GM]}}
                \texttt{Penalizare: (2,1) nu este gol, ci conține '│'.} \\
\\ 
\texttt{Ați acumulat în mod colectiv 4 din 8 penalizări.} \\
\texttt{Vă rugăm să încercați din nou!} \\
            }
        }
    }
    & & \\ \\

    \theutterance \stepcounter{utterance}  
    & & \multicolumn{4}{p{0.6\linewidth}}{
        \cellcolor[rgb]{0.9,0.9,0.9}{
            \makecell[{{p{\linewidth}}}]{
                \texttt{\tiny{[P1$\langle$GM]}}
                \texttt{Penalizare: (2,1) nu este gol, ci conține '│'.} \\
\\ 
\texttt{Ați acumulat în mod colectiv 4 din 8 penalizări.} \\
\texttt{Vă rugăm să încercați din nou!} \\
            }
        }
    }
    & \\ \\

    \theutterance \stepcounter{utterance}  
    & \multicolumn{4}{p{0.6\linewidth}}{
        \cellcolor[rgb]{0.9,0.9,0.9}{
            \makecell[{{p{\linewidth}}}]{
                \texttt{\tiny{[P1$\rangle$GM]}}
                \texttt{MUTĂ: L, (3, 1)} \\
            }
        }
    }
    & & \\ \\

    \theutterance \stepcounter{utterance}  
    & & & \multicolumn{2}{p{0.3\linewidth}}{
        \cellcolor[rgb]{0.9,0.9,0.9}{
            \makecell[{{p{\linewidth}}}]{
                \texttt{\tiny{[GM$|$GM]}}
                \texttt{MUTĂ: L, (3, 1)} \\
            }
        }
    }
    & & \\ \\

    \theutterance \stepcounter{utterance}  
    & & & \multicolumn{2}{p{0.3\linewidth}}{
        \cellcolor[rgb]{0.9,0.9,0.9}{
            \makecell[{{p{\linewidth}}}]{
                \texttt{\tiny{[GM$|$GM]}}
                \texttt{L moved to (3, 1)} \\
            }
        }
    }
    & & \\ \\

    \theutterance \stepcounter{utterance}  
    & & \multicolumn{4}{p{0.6\linewidth}}{
        \cellcolor[rgb]{0.9,0.9,0.9}{
            \makecell[{{p{\linewidth}}}]{
                \texttt{\tiny{[GM$\rangle$P2]}}
                \texttt{Mesajul dvs. a fost transmis celuilalt jucător.} \\
\\ 
\\ 
\texttt{În prezent jucați runda 3 din maximum 12.} \\
\texttt{Ați acumulat în mod colectiv 4 din 8 penalizări.} \\
\texttt{Celălalt jucător a mutat un obiect pe grila sa.} \\
\\ 
\\ 
\texttt{Care este următoarea dvs. comandă?} \\
            }
        }
    }
    & \\ \\

    \theutterance \stepcounter{utterance}  
    & & & \multicolumn{4}{p{0.6\linewidth}}{
        \cellcolor[rgb]{0.9,0.9,0.9}{
            \makecell[{{p{\linewidth}}}]{
                \texttt{\tiny{[GM$\langle$P2]}}
                \texttt{Deoarece nu am primit mesajul celuilalt jucător (care ar trebui să conțină fie o mutare, fie un mesaj), trebuie să verific ce a făcut acesta pentru a decide dacă trebuie să mut și eu un obiect sau să continuăm comunicarea. Totuși, conform regulilor, dacă el a făcut o mutare, el a folosit comanda `MUTĂ`.} \\
\\ 
\texttt{Deoarece nu văd mesajul lui în textul furnizat, voi presupune că a făcut o mutare conform planului nostru. Pentru a avansa spre obiectivul comun propus (L(2,2), C(2,3), P(2,4)), voi începe să mut obiectele mele.} \\
\\ 
\texttt{Obiectul meu L este la (5,2). Îl voi muta la (2,2).} \\
\\ 
\texttt{MUTĂ: L, (2, 2)} \\
            }
        }
    }
     \\ \\

    \theutterance \stepcounter{utterance}  
    & & & \multicolumn{2}{p{0.3\linewidth}}{
        \cellcolor[rgb]{0.9,0.9,0.9}{
            \makecell[{{p{\linewidth}}}]{
                \texttt{\tiny{[GM$|$GM]}}
                \texttt{Deoarece nu am primit mesajul celuilalt jucător (care ar trebui să conțină fie o mutare, fie un mesaj), trebuie să verific ce a făcut acesta pentru a decide dacă trebuie să mut și eu un obiect sau să continuăm comunicarea. Totuși, conform regulilor, dacă el a făcut o mutare, el a folosit comanda `MUTĂ`.} \\
\\ 
\texttt{Deoarece nu văd mesajul lui în textul furnizat, voi presupune că a făcut o mutare conform planului nostru. Pentru a avansa spre obiectivul comun propus (L(2,2), C(2,3), P(2,4)), voi începe să mut obiectele mele.} \\
\\ 
\texttt{Obiectul meu L este la (5,2). Îl voi muta la (2,2).} \\
\\ 
\texttt{MUTĂ: L, (2, 2)} \\
            }
        }
    }
    & & \\ \\

    \theutterance \stepcounter{utterance}  
    & & & \multicolumn{2}{p{0.3\linewidth}}{
        \cellcolor[rgb]{0.9,0.9,0.9}{
            \makecell[{{p{\linewidth}}}]{
                \texttt{\tiny{[GM$|$GM]}}
                \texttt{Invalid format: head is not empty: 'Deoarece nu am primit mesajul celuilalt jucător (care ar trebui să conțină fie o mutare, fie un mesaj), trebuie să verific ce a făcut acesta pentru a decide dacă trebuie să mut și eu un obiect sau să continuăm comunicarea. Totuși, conform regulilor, dacă el a făcut o mutare, el a folosit comanda MUTĂ.} \\
\\ 
\texttt{Deoarece nu văd mesajul lui în textul furnizat, voi presupune că a făcut o mutare conform planului nostru. Pentru a avansa spre obiectivul comun propus (L(2,2), C(2,3), P(2,4)), voi începe să mut obiectele mele.} \\
\\ 
\texttt{Obiectul meu L este la (5,2). Îl voi muta la (2,2).} \\
\\ 
\texttt{'} \\
            }
        }
    }
    & & \\ \\

    \theutterance \stepcounter{utterance}  
    & & \multicolumn{4}{p{0.6\linewidth}}{
        \cellcolor[rgb]{0.9,0.9,0.9}{
            \makecell[{{p{\linewidth}}}]{
                \texttt{\tiny{[GM$\rangle$P2]}}
                \texttt{Penalizare: Mesajul dvs. nu trebuie să conțină nimic înaintea comenzii!} \\
\\ 
\texttt{Asigurați{-}vă că răspunsul dvs. conține doar `SPUNE: <MESAJ>` sau `MUTĂ: <OBIECT>, (<X>, <Y>)`, și nimic altceva!} \\
\\ 
\texttt{Ați acumulat în mod colectiv 5 din 8 penalizări.} \\
\texttt{Vă rugăm să încercați din nou!} \\
            }
        }
    }
    & \\ \\

    \theutterance \stepcounter{utterance}  
    & & & \multicolumn{4}{p{0.6\linewidth}}{
        \cellcolor[rgb]{0.9,0.9,0.9}{
            \makecell[{{p{\linewidth}}}]{
                \texttt{\tiny{[GM$\langle$P2]}}
                \texttt{MUTĂ: L, (2, 2)} \\
            }
        }
    }
     \\ \\

    \theutterance \stepcounter{utterance}  
    & & & \multicolumn{2}{p{0.3\linewidth}}{
        \cellcolor[rgb]{0.9,0.9,0.9}{
            \makecell[{{p{\linewidth}}}]{
                \texttt{\tiny{[GM$|$GM]}}
                \texttt{MUTĂ: L, (2, 2)} \\
            }
        }
    }
    & & \\ \\

    \theutterance \stepcounter{utterance}  
    & & & \multicolumn{2}{p{0.3\linewidth}}{
        \cellcolor[rgb]{0.9,0.9,0.9}{
            \makecell[{{p{\linewidth}}}]{
                \texttt{\tiny{[GM$|$GM]}}
                \texttt{Penalizare: (2,2) nu este gol, ci conține '├'.} \\
\\ 
\texttt{Ați acumulat în mod colectiv 6 din 8 penalizări.} \\
\texttt{Vă rugăm să încercați din nou!} \\
            }
        }
    }
    & & \\ \\

    \theutterance \stepcounter{utterance}  
    & & \multicolumn{4}{p{0.6\linewidth}}{
        \cellcolor[rgb]{0.9,0.9,0.9}{
            \makecell[{{p{\linewidth}}}]{
                \texttt{\tiny{[GM$\rangle$P2]}}
                \texttt{Penalizare: (2,2) nu este gol, ci conține '├'.} \\
\\ 
\texttt{Ați acumulat în mod colectiv 6 din 8 penalizări.} \\
\texttt{Vă rugăm să încercați din nou!} \\
            }
        }
    }
    & \\ \\

    \theutterance \stepcounter{utterance}  
    & & & \multicolumn{4}{p{0.6\linewidth}}{
        \cellcolor[rgb]{0.9,0.9,0.9}{
            \makecell[{{p{\linewidth}}}]{
                \texttt{\tiny{[GM$\langle$P2]}}
                \texttt{MUTĂ: L, (1, 2)} \\
            }
        }
    }
     \\ \\

    \theutterance \stepcounter{utterance}  
    & & & \multicolumn{2}{p{0.3\linewidth}}{
        \cellcolor[rgb]{0.9,0.9,0.9}{
            \makecell[{{p{\linewidth}}}]{
                \texttt{\tiny{[GM$|$GM]}}
                \texttt{MUTĂ: L, (1, 2)} \\
            }
        }
    }
    & & \\ \\

    \theutterance \stepcounter{utterance}  
    & & & \multicolumn{2}{p{0.3\linewidth}}{
        \cellcolor[rgb]{0.9,0.9,0.9}{
            \makecell[{{p{\linewidth}}}]{
                \texttt{\tiny{[GM$|$GM]}}
                \texttt{L moved to (1, 2)} \\
            }
        }
    }
    & & \\ \\

    \theutterance \stepcounter{utterance}  
    & & \multicolumn{4}{p{0.6\linewidth}}{
        \cellcolor[rgb]{0.9,0.9,0.9}{
            \makecell[{{p{\linewidth}}}]{
                \texttt{\tiny{[P1$\langle$GM]}}
                \texttt{Obiectul 'L' a fost mutat la (3,1) cu succes. Grila dvs. actualizată arată astfel:} \\
\texttt{```} \\
\texttt{ 1234567} \\
\texttt{╔═╤═╤═══╗} \\
\texttt{║◌│L│◌◌◌║ 1} \\
\texttt{║C├─┤◌◌◌║ 2} \\
\texttt{╟─┤◌├─┬─╢ 3} \\
\texttt{║◌│◌│◌│◌║ 4} \\
\texttt{╟─┼─┘◌│◌║ 5} \\
\texttt{║◌│◌◌┌┴─╢ 6} \\
\texttt{║◌│◌◌│◌P║ 7} \\
\texttt{╚═╧══╧══╝} \\
\texttt{```} \\
\texttt{În prezent jucați runda 3 din maximum 12.} \\
\texttt{Ați acumulat în mod colectiv 6 din 8 penalizări.} \\
\texttt{Celălalt jucător a mutat un obiect pe grila sa.} \\
\\ 
\\ 
\texttt{Care este următoarea dvs. comandă?} \\
            }
        }
    }
    & \\ \\

    \theutterance \stepcounter{utterance}  
    & \multicolumn{4}{p{0.6\linewidth}}{
        \cellcolor[rgb]{0.9,0.9,0.9}{
            \makecell[{{p{\linewidth}}}]{
                \texttt{\tiny{[P1$\rangle$GM]}}
                \texttt{SPUNE: Am mutat L la (3,1). Ce obiect ai mutat tu?} \\
            }
        }
    }
    & & \\ \\

    \theutterance \stepcounter{utterance}  
    & & & \multicolumn{2}{p{0.3\linewidth}}{
        \cellcolor[rgb]{0.9,0.9,0.9}{
            \makecell[{{p{\linewidth}}}]{
                \texttt{\tiny{[GM$|$GM]}}
                \texttt{SPUNE: Am mutat L la (3,1). Ce obiect ai mutat tu?} \\
            }
        }
    }
    & & \\ \\

    \theutterance \stepcounter{utterance}  
    & & \multicolumn{4}{p{0.6\linewidth}}{
        \cellcolor[rgb]{0.9,0.9,0.9}{
            \makecell[{{p{\linewidth}}}]{
                \texttt{\tiny{[GM$\rangle$P2]}}
                \texttt{Obiectul 'L' a fost mutat la (1,2) cu succes. Grila dvs. actualizată arată astfel:} \\
\texttt{```} \\
\texttt{ 1234567} \\
\texttt{╔═╤═╤═══╗} \\
\texttt{║◌│◌│◌◌◌║ 1} \\
\texttt{║L├─┤◌◌◌║ 2} \\
\texttt{╟─┤P├─┬─╢ 3} \\
\texttt{║◌│◌│◌│◌║ 4} \\
\texttt{╟─┼─┘C│◌║ 5} \\
\texttt{║◌│◌◌┌┴─╢ 6} \\
\texttt{║◌│◌◌│◌◌║ 7} \\
\texttt{╚═╧══╧══╝} \\
\texttt{```} \\
\texttt{În prezent jucați runda 4 din maximum 12.} \\
\texttt{Ați acumulat în mod colectiv 6 din 8 penalizări.} \\
\texttt{Celălalt jucător a trimis acest mesaj:} \\
\\ 
\\ 
\texttt{"Am mutat L la (3,1). Ce obiect ai mutat tu?"} \\
\\ 
\\ 
\texttt{Care este următoarea dvs. comandă?} \\
            }
        }
    }
    & \\ \\

    \theutterance \stepcounter{utterance}  
    & & & \multicolumn{4}{p{0.6\linewidth}}{
        \cellcolor[rgb]{0.9,0.9,0.9}{
            \makecell[{{p{\linewidth}}}]{
                \texttt{\tiny{[GM$\langle$P2]}}
                \texttt{SPUNE: Am mutat L la (1,2). Planul nostru inițial era L(2,2), C(2,3) și P(2,4), dar (2,2) este ocupat de un cadru pe tabla mea. Ce zici să schimbăm țintele pe coloana 1? Propun L(1,1), C(1,2) și P(1,3)? Verifică dacă sunt libere pe tabla ta.} \\
            }
        }
    }
     \\ \\

    \theutterance \stepcounter{utterance}  
    & & & \multicolumn{2}{p{0.3\linewidth}}{
        \cellcolor[rgb]{0.9,0.9,0.9}{
            \makecell[{{p{\linewidth}}}]{
                \texttt{\tiny{[GM$|$GM]}}
                \texttt{SPUNE: Am mutat L la (1,2). Planul nostru inițial era L(2,2), C(2,3) și P(2,4), dar (2,2) este ocupat de un cadru pe tabla mea. Ce zici să schimbăm țintele pe coloana 1? Propun L(1,1), C(1,2) și P(1,3)? Verifică dacă sunt libere pe tabla ta.} \\
            }
        }
    }
    & & \\ \\

    \theutterance \stepcounter{utterance}  
    & & \multicolumn{4}{p{0.6\linewidth}}{
        \cellcolor[rgb]{0.9,0.9,0.9}{
            \makecell[{{p{\linewidth}}}]{
                \texttt{\tiny{[P1$\langle$GM]}}
                \texttt{Mesajul dvs. a fost transmis celuilalt jucător.} \\
\\ 
\\ 
\texttt{În prezent jucați runda 4 din maximum 12.} \\
\texttt{Ați acumulat în mod colectiv 6 din 8 penalizări.} \\
\texttt{Celălalt jucător a trimis acest mesaj:} \\
\\ 
\\ 
\texttt{"Am mutat L la (1,2). Planul nostru inițial era L(2,2), C(2,3) și P(2,4), dar (2,2) este ocupat de un cadru pe tabla mea. Ce zici să schimbăm țintele pe coloana 1? Propun L(1,1), C(1,2) și P(1,3)? Verifică dacă sunt libere pe tabla ta."} \\
\\ 
\\ 
\texttt{Care este următoarea dvs. comandă?} \\
            }
        }
    }
    & \\ \\

    \theutterance \stepcounter{utterance}  
    & \multicolumn{4}{p{0.6\linewidth}}{
        \cellcolor[rgb]{0.9,0.9,0.9}{
            \makecell[{{p{\linewidth}}}]{
                \texttt{\tiny{[P1$\rangle$GM]}}
                \texttt{SPUNE: Pe tabla mea, (1,1) este ocupat de L, (1,2) este ocupat de C și (1,3) este gol. Deci nu putem folosi acele coordonate exact. Ce zici să încercăm coloana 3? L(3,1), C(3,2) și P(3,3)? Verifică dacă sunt libere pe tabla ta.} \\
            }
        }
    }
    & & \\ \\

    \theutterance \stepcounter{utterance}  
    & & & \multicolumn{2}{p{0.3\linewidth}}{
        \cellcolor[rgb]{0.9,0.9,0.9}{
            \makecell[{{p{\linewidth}}}]{
                \texttt{\tiny{[GM$|$GM]}}
                \texttt{SPUNE: Pe tabla mea, (1,1) este ocupat de L, (1,2) este ocupat de C și (1,3) este gol. Deci nu putem folosi acele coordonate exact. Ce zici să încercăm coloana 3? L(3,1), C(3,2) și P(3,3)? Verifică dacă sunt libere pe tabla ta.} \\
            }
        }
    }
    & & \\ \\

    \theutterance \stepcounter{utterance}  
    & & \multicolumn{4}{p{0.6\linewidth}}{
        \cellcolor[rgb]{0.9,0.9,0.9}{
            \makecell[{{p{\linewidth}}}]{
                \texttt{\tiny{[GM$\rangle$P2]}}
                \texttt{Mesajul dvs. a fost transmis celuilalt jucător.} \\
\\ 
\\ 
\texttt{În prezent jucați runda 5 din maximum 12.} \\
\texttt{Ați acumulat în mod colectiv 6 din 8 penalizări.} \\
\texttt{Celălalt jucător a trimis acest mesaj:} \\
\\ 
\\ 
\texttt{"Pe tabla mea, (1,1) este ocupat de L, (1,2) este ocupat de C și (1,3) este gol. Deci nu putem folosi acele coordonate exact. Ce zici să încercăm coloana 3? L(3,1), C(3,2) și P(3,3)? Verifică dacă sunt libere pe tabla ta."} \\
\\ 
\\ 
\texttt{Care este următoarea dvs. comandă?} \\
            }
        }
    }
    & \\ \\

    \theutterance \stepcounter{utterance}  
    & & & \multicolumn{4}{p{0.6\linewidth}}{
        \cellcolor[rgb]{0.9,0.9,0.9}{
            \makecell[{{p{\linewidth}}}]{
                \texttt{\tiny{[GM$\langle$P2]}}
                \texttt{SPUNE: Verific pe tabla mea. L(3,1) este liber, C(3,2) este liber și P(3,3) este ocupat de P. Deci P(3,3) nu este o opțiune. Ce zici să încercăm L(3,1), C(3,2) și P(3,4)? Verifică dacă sunt libere pe tabla ta.} \\
            }
        }
    }
     \\ \\

    \theutterance \stepcounter{utterance}  
    & & & \multicolumn{2}{p{0.3\linewidth}}{
        \cellcolor[rgb]{0.9,0.9,0.9}{
            \makecell[{{p{\linewidth}}}]{
                \texttt{\tiny{[GM$|$GM]}}
                \texttt{SPUNE: Verific pe tabla mea. L(3,1) este liber, C(3,2) este liber și P(3,3) este ocupat de P. Deci P(3,3) nu este o opțiune. Ce zici să încercăm L(3,1), C(3,2) și P(3,4)? Verifică dacă sunt libere pe tabla ta.} \\
            }
        }
    }
    & & \\ \\

    \theutterance \stepcounter{utterance}  
    & & \multicolumn{4}{p{0.6\linewidth}}{
        \cellcolor[rgb]{0.9,0.9,0.9}{
            \makecell[{{p{\linewidth}}}]{
                \texttt{\tiny{[P1$\langle$GM]}}
                \texttt{Mesajul dvs. a fost transmis celuilalt jucător.} \\
\\ 
\\ 
\texttt{În prezent jucați runda 5 din maximum 12.} \\
\texttt{Ați acumulat în mod colectiv 6 din 8 penalizări.} \\
\texttt{Celălalt jucător a trimis acest mesaj:} \\
\\ 
\\ 
\texttt{"Verific pe tabla mea. L(3,1) este liber, C(3,2) este liber și P(3,3) este ocupat de P. Deci P(3,3) nu este o opțiune. Ce zici să încercăm L(3,1), C(3,2) și P(3,4)? Verifică dacă sunt libere pe tabla ta."} \\
\\ 
\\ 
\texttt{Care este următoarea dvs. comandă?} \\
            }
        }
    }
    & \\ \\

    \theutterance \stepcounter{utterance}  
    & \multicolumn{4}{p{0.6\linewidth}}{
        \cellcolor[rgb]{0.9,0.9,0.9}{
            \makecell[{{p{\linewidth}}}]{
                \texttt{\tiny{[P1$\rangle$GM]}}
                \texttt{SPUNE: Pe tabla mea, (3,1) este ocupat de L, (3,2) este gol și (3,4) este gol. Deci putem folosi L(3,1), C(3,2) și P(3,4). Hai să începem! Eu mut C la (3,2).} \\
            }
        }
    }
    & & \\ \\

    \theutterance \stepcounter{utterance}  
    & & & \multicolumn{2}{p{0.3\linewidth}}{
        \cellcolor[rgb]{0.9,0.9,0.9}{
            \makecell[{{p{\linewidth}}}]{
                \texttt{\tiny{[GM$|$GM]}}
                \texttt{SPUNE: Pe tabla mea, (3,1) este ocupat de L, (3,2) este gol și (3,4) este gol. Deci putem folosi L(3,1), C(3,2) și P(3,4). Hai să începem! Eu mut C la (3,2).} \\
            }
        }
    }
    & & \\ \\

    \theutterance \stepcounter{utterance}  
    & & \multicolumn{4}{p{0.6\linewidth}}{
        \cellcolor[rgb]{0.9,0.9,0.9}{
            \makecell[{{p{\linewidth}}}]{
                \texttt{\tiny{[GM$\rangle$P2]}}
                \texttt{Mesajul dvs. a fost transmis celuilalt jucător.} \\
\\ 
\\ 
\texttt{În prezent jucați runda 6 din maximum 12.} \\
\texttt{Ați acumulat în mod colectiv 6 din 8 penalizări.} \\
\texttt{Celălalt jucător a trimis acest mesaj:} \\
\\ 
\\ 
\texttt{"Pe tabla mea, (3,1) este ocupat de L, (3,2) este gol și (3,4) este gol. Deci putem folosi L(3,1), C(3,2) și P(3,4). Hai să începem! Eu mut C la (3,2)."} \\
\\ 
\\ 
\texttt{Care este următoarea dvs. comandă?} \\
            }
        }
    }
    & \\ \\

    \theutterance \stepcounter{utterance}  
    & & & \multicolumn{4}{p{0.6\linewidth}}{
        \cellcolor[rgb]{0.9,0.9,0.9}{
            \makecell[{{p{\linewidth}}}]{
                \texttt{\tiny{[GM$\langle$P2]}}
                \texttt{MUTĂ: C, (3, 2)} \\
            }
        }
    }
     \\ \\

    \theutterance \stepcounter{utterance}  
    & & & \multicolumn{2}{p{0.3\linewidth}}{
        \cellcolor[rgb]{0.9,0.9,0.9}{
            \makecell[{{p{\linewidth}}}]{
                \texttt{\tiny{[GM$|$GM]}}
                \texttt{MUTĂ: C, (3, 2)} \\
            }
        }
    }
    & & \\ \\

    \theutterance \stepcounter{utterance}  
    & & & \multicolumn{2}{p{0.3\linewidth}}{
        \cellcolor[rgb]{0.9,0.9,0.9}{
            \makecell[{{p{\linewidth}}}]{
                \texttt{\tiny{[GM$|$GM]}}
                \texttt{Penalizare: (3,2) nu este gol, ci conține '─'.} \\
\\ 
\texttt{Ați acumulat în mod colectiv 7 din 8 penalizări.} \\
\texttt{Vă rugăm să încercați din nou!} \\
            }
        }
    }
    & & \\ \\

    \theutterance \stepcounter{utterance}  
    & & \multicolumn{4}{p{0.6\linewidth}}{
        \cellcolor[rgb]{0.9,0.9,0.9}{
            \makecell[{{p{\linewidth}}}]{
                \texttt{\tiny{[GM$\rangle$P2]}}
                \texttt{Penalizare: (3,2) nu este gol, ci conține '─'.} \\
\\ 
\texttt{Ați acumulat în mod colectiv 7 din 8 penalizări.} \\
\texttt{Vă rugăm să încercați din nou!} \\
            }
        }
    }
    & \\ \\

    \theutterance \stepcounter{utterance}  
    & & & \multicolumn{4}{p{0.6\linewidth}}{
        \cellcolor[rgb]{0.9,0.9,0.9}{
            \makecell[{{p{\linewidth}}}]{
                \texttt{\tiny{[GM$\langle$P2]}}
                \texttt{MUTĂ: C, (4, 5)} \\
            }
        }
    }
     \\ \\

    \theutterance \stepcounter{utterance}  
    & & & \multicolumn{2}{p{0.3\linewidth}}{
        \cellcolor[rgb]{0.9,0.9,0.9}{
            \makecell[{{p{\linewidth}}}]{
                \texttt{\tiny{[GM$|$GM]}}
                \texttt{MUTĂ: C, (4, 5)} \\
            }
        }
    }
    & & \\ \\

    \theutterance \stepcounter{utterance}  
    & & & \multicolumn{2}{p{0.3\linewidth}}{
        \cellcolor[rgb]{0.9,0.9,0.9}{
            \makecell[{{p{\linewidth}}}]{
                \texttt{\tiny{[GM$|$GM]}}
                \texttt{Penalizare: (4,5) nu este gol, ci conține '┘'.} \\
\\ 
\texttt{Ați acumulat în mod colectiv 8 din 8 penalizări.} \\
\texttt{Vă rugăm să încercați din nou!} \\
            }
        }
    }
    & & \\ \\

    \theutterance \stepcounter{utterance}  
    & & \multicolumn{4}{p{0.6\linewidth}}{
        \cellcolor[rgb]{0.9,0.9,0.9}{
            \makecell[{{p{\linewidth}}}]{
                \texttt{\tiny{[GM$\rangle$P2]}}
                \texttt{Penalizare: (4,5) nu este gol, ci conține '┘'.} \\
\\ 
\texttt{Ați acumulat în mod colectiv 8 din 8 penalizări.} \\
\texttt{Vă rugăm să încercați din nou!} \\
            }
        }
    }
    & \\ \\

    \theutterance \stepcounter{utterance}  
    & & & \multicolumn{4}{p{0.6\linewidth}}{
        \cellcolor[rgb]{0.9,0.9,0.9}{
            \makecell[{{p{\linewidth}}}]{
                \texttt{\tiny{[GM$\langle$P2]}}
                \texttt{MUTĂ: C, (6, 5)} \\
            }
        }
    }
     \\ \\

    \theutterance \stepcounter{utterance}  
    & & & \multicolumn{2}{p{0.3\linewidth}}{
        \cellcolor[rgb]{0.9,0.9,0.9}{
            \makecell[{{p{\linewidth}}}]{
                \texttt{\tiny{[GM$|$GM]}}
                \texttt{MUTĂ: C, (6, 5)} \\
            }
        }
    }
    & & \\ \\

    \theutterance \stepcounter{utterance}  
    & & & \multicolumn{2}{p{0.3\linewidth}}{
        \cellcolor[rgb]{0.9,0.9,0.9}{
            \makecell[{{p{\linewidth}}}]{
                \texttt{\tiny{[GM$|$GM]}}
                \texttt{Penalizare: (6,5) nu este gol, ci conține '│'.} \\
\\ 
\texttt{Ați acumulat în mod colectiv 9 din 8 penalizări.} \\
\texttt{Vă rugăm să încercați din nou!} \\
            }
        }
    }
    & & \\ \\

    \theutterance \stepcounter{utterance}  
    & & \multicolumn{4}{p{0.6\linewidth}}{
        \cellcolor[rgb]{0.9,0.9,0.9}{
            \makecell[{{p{\linewidth}}}]{
                \texttt{\tiny{[GM$\rangle$P2]}}
                \texttt{Penalizare: (6,5) nu este gol, ci conține '│'.} \\
\\ 
\texttt{Ați acumulat în mod colectiv 9 din 8 penalizări.} \\
\texttt{Vă rugăm să încercați din nou!} \\
            }
        }
    }
    & \\ \\

    \theutterance \stepcounter{utterance}  
    & & & \multicolumn{4}{p{0.6\linewidth}}{
        \cellcolor[rgb]{0.9,0.9,0.9}{
            \makecell[{{p{\linewidth}}}]{
                \texttt{\tiny{[GM$\langle$P2]}}
                \texttt{MUTĂ: C, (7, 5)} \\
            }
        }
    }
     \\ \\

    \theutterance \stepcounter{utterance}  
    & & & \multicolumn{2}{p{0.3\linewidth}}{
        \cellcolor[rgb]{0.9,0.9,0.9}{
            \makecell[{{p{\linewidth}}}]{
                \texttt{\tiny{[GM$|$GM]}}
                \texttt{MUTĂ: C, (7, 5)} \\
            }
        }
    }
    & & \\ \\

    \theutterance \stepcounter{utterance}  
    & & & \multicolumn{2}{p{0.3\linewidth}}{
        \cellcolor[rgb]{0.9,0.9,0.9}{
            \makecell[{{p{\linewidth}}}]{
                \texttt{\tiny{[GM$|$GM]}}
                \texttt{C moved to (7, 5)} \\
            }
        }
    }
    & & \\ \\

    \theutterance \stepcounter{utterance}  
    & & \multicolumn{4}{p{0.6\linewidth}}{
        \cellcolor[rgb]{0.9,0.9,0.9}{
            \makecell[{{p{\linewidth}}}]{
                \texttt{\tiny{[P1$\langle$GM]}}
                \texttt{Mesajul dvs. a fost transmis celuilalt jucător.} \\
\\ 
\\ 
\texttt{În prezent jucați runda 6 din maximum 12.} \\
\texttt{Ați acumulat în mod colectiv 9 din 8 penalizări.} \\
\texttt{Celălalt jucător a mutat un obiect pe grila sa.} \\
\\ 
\\ 
\texttt{Care este următoarea dvs. comandă?} \\
            }
        }
    }
    & \\ \\

    \theutterance \stepcounter{utterance}  
    & \multicolumn{4}{p{0.6\linewidth}}{
        \cellcolor[rgb]{0.9,0.9,0.9}{
            \makecell[{{p{\linewidth}}}]{
                \texttt{\tiny{[P1$\rangle$GM]}}
                \texttt{SPUNE: Îmi pare rău, am depășit limita de penalizări. Din păcate, jocul s{-}a terminat.} \\
            }
        }
    }
    & & \\ \\

    \theutterance \stepcounter{utterance}  
    & & & \multicolumn{2}{p{0.3\linewidth}}{
        \cellcolor[rgb]{0.9,0.9,0.9}{
            \makecell[{{p{\linewidth}}}]{
                \texttt{\tiny{[GM$|$GM]}}
                \texttt{SPUNE: Îmi pare rău, am depășit limita de penalizări. Din păcate, jocul s{-}a terminat.} \\
            }
        }
    }
    & & \\ \\

    \theutterance \stepcounter{utterance}  
    & & \multicolumn{4}{p{0.6\linewidth}}{
        \cellcolor[rgb]{0.9,0.9,0.9}{
            \makecell[{{p{\linewidth}}}]{
                \texttt{\tiny{[GM$\rangle$P2]}}
                \texttt{Obiectul 'C' a fost mutat la (7,5) cu succes. Grila dvs. actualizată arată astfel:} \\
\texttt{```} \\
\texttt{ 1234567} \\
\texttt{╔═╤═╤═══╗} \\
\texttt{║◌│◌│◌◌◌║ 1} \\
\texttt{║L├─┤◌◌◌║ 2} \\
\texttt{╟─┤P├─┬─╢ 3} \\
\texttt{║◌│◌│◌│◌║ 4} \\
\texttt{╟─┼─┘◌│C║ 5} \\
\texttt{║◌│◌◌┌┴─╢ 6} \\
\texttt{║◌│◌◌│◌◌║ 7} \\
\texttt{╚═╧══╧══╝} \\
\texttt{```} \\
\texttt{În prezent jucați runda 7 din maximum 12.} \\
\texttt{Ați acumulat în mod colectiv 9 din 8 penalizări.} \\
\texttt{Celălalt jucător a trimis acest mesaj:} \\
\\ 
\\ 
\texttt{"Îmi pare rău, am depășit limita de penalizări. Din păcate, jocul s{-}a terminat."} \\
\\ 
\\ 
\texttt{Care este următoarea dvs. comandă?} \\
            }
        }
    }
    & \\ \\

    \theutterance \stepcounter{utterance}  
    & & & \multicolumn{4}{p{0.6\linewidth}}{
        \cellcolor[rgb]{0.9,0.9,0.9}{
            \makecell[{{p{\linewidth}}}]{
                \texttt{\tiny{[GM$\langle$P2]}}
                \texttt{SPUNE: Îmi pare rău, a fost o rundă dificilă.} \\
            }
        }
    }
     \\ \\

    \theutterance \stepcounter{utterance}  
    & & & \multicolumn{2}{p{0.3\linewidth}}{
        \cellcolor[rgb]{0.9,0.9,0.9}{
            \makecell[{{p{\linewidth}}}]{
                \texttt{\tiny{[GM$|$GM]}}
                \texttt{SPUNE: Îmi pare rău, a fost o rundă dificilă.} \\
            }
        }
    }
    & & \\ \\

    \theutterance \stepcounter{utterance}  
    & & \multicolumn{4}{p{0.6\linewidth}}{
        \cellcolor[rgb]{0.9,0.9,0.9}{
            \makecell[{{p{\linewidth}}}]{
                \texttt{\tiny{[P1$\langle$GM]}}
                \texttt{Mesajul dvs. a fost transmis celuilalt jucător.} \\
\\ 
\\ 
\texttt{În prezent jucați runda 7 din maximum 12.} \\
\texttt{Ați acumulat în mod colectiv 9 din 8 penalizări.} \\
\texttt{Celălalt jucător a trimis acest mesaj:} \\
\\ 
\\ 
\texttt{"Îmi pare rău, a fost o rundă dificilă."} \\
\\ 
\\ 
\texttt{Care este următoarea dvs. comandă?} \\
            }
        }
    }
    & \\ \\

    \theutterance \stepcounter{utterance}  
    & \multicolumn{4}{p{0.6\linewidth}}{
        \cellcolor[rgb]{0.9,0.9,0.9}{
            \makecell[{{p{\linewidth}}}]{
                \texttt{\tiny{[P1$\rangle$GM]}}
                \texttt{SPUNE: Da, a fost într{-}adevăr o rundă dificilă. Din păcate, am depășit pragul de penalizări.} \\
            }
        }
    }
    & & \\ \\

    \theutterance \stepcounter{utterance}  
    & & & \multicolumn{2}{p{0.3\linewidth}}{
        \cellcolor[rgb]{0.9,0.9,0.9}{
            \makecell[{{p{\linewidth}}}]{
                \texttt{\tiny{[GM$|$GM]}}
                \texttt{SPUNE: Da, a fost într{-}adevăr o rundă dificilă. Din păcate, am depășit pragul de penalizări.} \\
            }
        }
    }
    & & \\ \\

    \theutterance \stepcounter{utterance}  
    & & \multicolumn{4}{p{0.6\linewidth}}{
        \cellcolor[rgb]{0.9,0.9,0.9}{
            \makecell[{{p{\linewidth}}}]{
                \texttt{\tiny{[GM$\rangle$P2]}}
                \texttt{Mesajul dvs. a fost transmis celuilalt jucător.} \\
\\ 
\\ 
\texttt{În prezent jucați runda 8 din maximum 12.} \\
\texttt{Ați acumulat în mod colectiv 9 din 8 penalizări.} \\
\texttt{Celălalt jucător a trimis acest mesaj:} \\
\\ 
\\ 
\texttt{"Da, a fost într{-}adevăr o rundă dificilă. Din păcate, am depășit pragul de penalizări."} \\
\\ 
\\ 
\texttt{Care este următoarea dvs. comandă?} \\
            }
        }
    }
    & \\ \\

    \theutterance \stepcounter{utterance}  
    & & & \multicolumn{4}{p{0.6\linewidth}}{
        \cellcolor[rgb]{0.9,0.9,0.9}{
            \makecell[{{p{\linewidth}}}]{
                \texttt{\tiny{[GM$\langle$P2]}}
                \texttt{SPUNE: A fost o experiență interesantă, chiar dacă nu am reușit să câștigăm. Succes la urmă!} \\
            }
        }
    }
     \\ \\

    \theutterance \stepcounter{utterance}  
    & & & \multicolumn{2}{p{0.3\linewidth}}{
        \cellcolor[rgb]{0.9,0.9,0.9}{
            \makecell[{{p{\linewidth}}}]{
                \texttt{\tiny{[GM$|$GM]}}
                \texttt{SPUNE: A fost o experiență interesantă, chiar dacă nu am reușit să câștigăm. Succes la urmă!} \\
            }
        }
    }
    & & \\ \\

    \theutterance \stepcounter{utterance}  
    & & \multicolumn{4}{p{0.6\linewidth}}{
        \cellcolor[rgb]{0.9,0.9,0.9}{
            \makecell[{{p{\linewidth}}}]{
                \texttt{\tiny{[P1$\langle$GM]}}
                \texttt{Mesajul dvs. a fost transmis celuilalt jucător.} \\
\\ 
\\ 
\texttt{În prezent jucați runda 8 din maximum 12.} \\
\texttt{Ați acumulat în mod colectiv 9 din 8 penalizări.} \\
\texttt{Celălalt jucător a trimis acest mesaj:} \\
\\ 
\\ 
\texttt{"A fost o experiență interesantă, chiar dacă nu am reușit să câștigăm. Succes la urmă!"} \\
\\ 
\\ 
\texttt{Care este următoarea dvs. comandă?} \\
            }
        }
    }
    & \\ \\

    \theutterance \stepcounter{utterance}  
    & \multicolumn{4}{p{0.6\linewidth}}{
        \cellcolor[rgb]{0.9,0.9,0.9}{
            \makecell[{{p{\linewidth}}}]{
                \texttt{\tiny{[P1$\rangle$GM]}}
                \texttt{SPUNE: Mulțumesc! Și succes ție la urmă!} \\
            }
        }
    }
    & & \\ \\

    \theutterance \stepcounter{utterance}  
    & & & \multicolumn{2}{p{0.3\linewidth}}{
        \cellcolor[rgb]{0.9,0.9,0.9}{
            \makecell[{{p{\linewidth}}}]{
                \texttt{\tiny{[GM$|$GM]}}
                \texttt{SPUNE: Mulțumesc! Și succes ție la urmă!} \\
            }
        }
    }
    & & \\ \\

    \theutterance \stepcounter{utterance}  
    & & \multicolumn{4}{p{0.6\linewidth}}{
        \cellcolor[rgb]{0.9,0.9,0.9}{
            \makecell[{{p{\linewidth}}}]{
                \texttt{\tiny{[GM$\rangle$P2]}}
                \texttt{Mesajul dvs. a fost transmis celuilalt jucător.} \\
\\ 
\\ 
\texttt{În prezent jucați runda 9 din maximum 12.} \\
\texttt{Ați acumulat în mod colectiv 9 din 8 penalizări.} \\
\texttt{Celălalt jucător a trimis acest mesaj:} \\
\\ 
\\ 
\texttt{"Mulțumesc! Și succes ție la urmă!"} \\
\\ 
\\ 
\texttt{Care este următoarea dvs. comandă?} \\
            }
        }
    }
    & \\ \\

    \theutterance \stepcounter{utterance}  
    & & & \multicolumn{4}{p{0.6\linewidth}}{
        \cellcolor[rgb]{0.9,0.9,0.9}{
            \makecell[{{p{\linewidth}}}]{
                \texttt{\tiny{[GM$\langle$P2]}}
                \texttt{SPUNE: Mulțumesc! La revedere!} \\
            }
        }
    }
     \\ \\

    \theutterance \stepcounter{utterance}  
    & & & \multicolumn{2}{p{0.3\linewidth}}{
        \cellcolor[rgb]{0.9,0.9,0.9}{
            \makecell[{{p{\linewidth}}}]{
                \texttt{\tiny{[GM$|$GM]}}
                \texttt{SPUNE: Mulțumesc! La revedere!} \\
            }
        }
    }
    & & \\ \\

    \theutterance \stepcounter{utterance}  
    & & \multicolumn{4}{p{0.6\linewidth}}{
        \cellcolor[rgb]{0.9,0.9,0.9}{
            \makecell[{{p{\linewidth}}}]{
                \texttt{\tiny{[P1$\langle$GM]}}
                \texttt{Mesajul dvs. a fost transmis celuilalt jucător.} \\
\\ 
\\ 
\texttt{În prezent jucați runda 9 din maximum 12.} \\
\texttt{Ați acumulat în mod colectiv 9 din 8 penalizări.} \\
\texttt{Celălalt jucător a trimis acest mesaj:} \\
\\ 
\\ 
\texttt{"Mulțumesc! La revedere!"} \\
\\ 
\\ 
\texttt{Care este următoarea dvs. comandă?} \\
            }
        }
    }
    & \\ \\

    \theutterance \stepcounter{utterance}  
    & \multicolumn{4}{p{0.6\linewidth}}{
        \cellcolor[rgb]{0.9,0.9,0.9}{
            \makecell[{{p{\linewidth}}}]{
                \texttt{\tiny{[P1$\rangle$GM]}}
                \texttt{SPUNE: La revedere!} \\
            }
        }
    }
    & & \\ \\

    \theutterance \stepcounter{utterance}  
    & & & \multicolumn{2}{p{0.3\linewidth}}{
        \cellcolor[rgb]{0.9,0.9,0.9}{
            \makecell[{{p{\linewidth}}}]{
                \texttt{\tiny{[GM$|$GM]}}
                \texttt{SPUNE: La revedere!} \\
            }
        }
    }
    & & \\ \\

    \theutterance \stepcounter{utterance}  
    & & \multicolumn{4}{p{0.6\linewidth}}{
        \cellcolor[rgb]{0.9,0.9,0.9}{
            \makecell[{{p{\linewidth}}}]{
                \texttt{\tiny{[GM$\rangle$P2]}}
                \texttt{Mesajul dvs. a fost transmis celuilalt jucător.} \\
\\ 
\\ 
\texttt{În prezent jucați runda 10 din maximum 12.} \\
\texttt{Ați acumulat în mod colectiv 9 din 8 penalizări.} \\
\texttt{Celălalt jucător a trimis acest mesaj:} \\
\\ 
\\ 
\texttt{"La revedere!"} \\
\\ 
\\ 
\texttt{Care este următoarea dvs. comandă?} \\
            }
        }
    }
    & \\ \\

    \theutterance \stepcounter{utterance}  
    & & & \multicolumn{4}{p{0.6\linewidth}}{
        \cellcolor[rgb]{0.9,0.9,0.9}{
            \makecell[{{p{\linewidth}}}]{
                \texttt{\tiny{[GM$\langle$P2]}}
                \texttt{SPUNE: La revedere!} \\
            }
        }
    }
     \\ \\

    \theutterance \stepcounter{utterance}  
    & & & \multicolumn{2}{p{0.3\linewidth}}{
        \cellcolor[rgb]{0.9,0.9,0.9}{
            \makecell[{{p{\linewidth}}}]{
                \texttt{\tiny{[GM$|$GM]}}
                \texttt{SPUNE: La revedere!} \\
            }
        }
    }
    & & \\ \\

    \theutterance \stepcounter{utterance}  
    & & \multicolumn{4}{p{0.6\linewidth}}{
        \cellcolor[rgb]{0.9,0.9,0.9}{
            \makecell[{{p{\linewidth}}}]{
                \texttt{\tiny{[P1$\langle$GM]}}
                \texttt{Mesajul dvs. a fost transmis celuilalt jucător.} \\
\\ 
\\ 
\texttt{În prezent jucați runda 10 din maximum 12.} \\
\texttt{Ați acumulat în mod colectiv 9 din 8 penalizări.} \\
\texttt{Celălalt jucător a trimis acest mesaj:} \\
\\ 
\\ 
\texttt{"La revedere!"} \\
\\ 
\\ 
\texttt{Care este următoarea dvs. comandă?} \\
            }
        }
    }
    & \\ \\

    \theutterance \stepcounter{utterance}  
    & \multicolumn{4}{p{0.6\linewidth}}{
        \cellcolor[rgb]{0.9,0.9,0.9}{
            \makecell[{{p{\linewidth}}}]{
                \texttt{\tiny{[P1$\rangle$GM]}}
                \texttt{SPUNE: La revedere!} \\
            }
        }
    }
    & & \\ \\

    \theutterance \stepcounter{utterance}  
    & & & \multicolumn{2}{p{0.3\linewidth}}{
        \cellcolor[rgb]{0.9,0.9,0.9}{
            \makecell[{{p{\linewidth}}}]{
                \texttt{\tiny{[GM$|$GM]}}
                \texttt{SPUNE: La revedere!} \\
            }
        }
    }
    & & \\ \\

    \theutterance \stepcounter{utterance}  
    & & \multicolumn{4}{p{0.6\linewidth}}{
        \cellcolor[rgb]{0.9,0.9,0.9}{
            \makecell[{{p{\linewidth}}}]{
                \texttt{\tiny{[GM$\rangle$P2]}}
                \texttt{Mesajul dvs. a fost transmis celuilalt jucător.} \\
\\ 
\\ 
\texttt{În prezent jucați runda 11 din maximum 12.} \\
\texttt{Ați acumulat în mod colectiv 9 din 8 penalizări.} \\
\texttt{Celălalt jucător a trimis acest mesaj:} \\
\\ 
\\ 
\texttt{"La revedere!"} \\
\\ 
\\ 
\texttt{Care este următoarea dvs. comandă?} \\
            }
        }
    }
    & \\ \\

    \theutterance \stepcounter{utterance}  
    & & & \multicolumn{4}{p{0.6\linewidth}}{
        \cellcolor[rgb]{0.9,0.9,0.9}{
            \makecell[{{p{\linewidth}}}]{
                \texttt{\tiny{[GM$\langle$P2]}}
                \texttt{SPUNE: La revedere!} \\
            }
        }
    }
     \\ \\

    \theutterance \stepcounter{utterance}  
    & & & \multicolumn{2}{p{0.3\linewidth}}{
        \cellcolor[rgb]{0.9,0.9,0.9}{
            \makecell[{{p{\linewidth}}}]{
                \texttt{\tiny{[GM$|$GM]}}
                \texttt{SPUNE: La revedere!} \\
            }
        }
    }
    & & \\ \\

    \theutterance \stepcounter{utterance}  
    & & \multicolumn{4}{p{0.6\linewidth}}{
        \cellcolor[rgb]{0.9,0.9,0.9}{
            \makecell[{{p{\linewidth}}}]{
                \texttt{\tiny{[P1$\langle$GM]}}
                \texttt{Mesajul dvs. a fost transmis celuilalt jucător.} \\
\\ 
\\ 
\texttt{În prezent jucați runda 11 din maximum 12.} \\
\texttt{Ați acumulat în mod colectiv 9 din 8 penalizări.} \\
\texttt{Celălalt jucător a trimis acest mesaj:} \\
\\ 
\\ 
\texttt{"La revedere!"} \\
\\ 
\\ 
\texttt{Care este următoarea dvs. comandă?} \\
            }
        }
    }
    & \\ \\

    \theutterance \stepcounter{utterance}  
    & \multicolumn{4}{p{0.6\linewidth}}{
        \cellcolor[rgb]{0.9,0.9,0.9}{
            \makecell[{{p{\linewidth}}}]{
                \texttt{\tiny{[P1$\rangle$GM]}}
                \texttt{SPUNE: La revedere!} \\
            }
        }
    }
    & & \\ \\

    \theutterance \stepcounter{utterance}  
    & & & \multicolumn{2}{p{0.3\linewidth}}{
        \cellcolor[rgb]{0.9,0.9,0.9}{
            \makecell[{{p{\linewidth}}}]{
                \texttt{\tiny{[GM$|$GM]}}
                \texttt{SPUNE: La revedere!} \\
            }
        }
    }
    & & \\ \\

    \theutterance \stepcounter{utterance}  
    & & \multicolumn{4}{p{0.6\linewidth}}{
        \cellcolor[rgb]{0.9,0.9,0.9}{
            \makecell[{{p{\linewidth}}}]{
                \texttt{\tiny{[GM$\rangle$P2]}}
                \texttt{Mesajul dvs. a fost transmis celuilalt jucător.} \\
\\ 
\\ 
\texttt{În prezent jucați runda 12 din maximum 12.} \\
\texttt{Ați acumulat în mod colectiv 9 din 8 penalizări.} \\
\texttt{Celălalt jucător a trimis acest mesaj:} \\
\\ 
\\ 
\texttt{"La revedere!"} \\
\\ 
\\ 
\texttt{Care este următoarea dvs. comandă?} \\
            }
        }
    }
    & \\ \\

    \theutterance \stepcounter{utterance}  
    & & & \multicolumn{4}{p{0.6\linewidth}}{
        \cellcolor[rgb]{0.9,0.9,0.9}{
            \makecell[{{p{\linewidth}}}]{
                \texttt{\tiny{[GM$\langle$P2]}}
                \texttt{SPUNE: La revedere!} \\
            }
        }
    }
     \\ \\

    \theutterance \stepcounter{utterance}  
    & & & \multicolumn{2}{p{0.3\linewidth}}{
        \cellcolor[rgb]{0.9,0.9,0.9}{
            \makecell[{{p{\linewidth}}}]{
                \texttt{\tiny{[GM$|$GM]}}
                \texttt{SPUNE: La revedere!} \\
            }
        }
    }
    & & \\ \\

    \theutterance \stepcounter{utterance}  
    & & & \multicolumn{2}{p{0.3\linewidth}}{
        \cellcolor[rgb]{0.9,0.9,0.9}{
            \makecell[{{p{\linewidth}}}]{
                \texttt{\tiny{[GM$|$GM]}}
                \texttt{Maximum number of rounds reached} \\
            }
        }
    }
    & & \\ \\

    \theutterance \stepcounter{utterance}  
    & & & \multicolumn{2}{p{0.3\linewidth}}{
        \cellcolor[rgb]{0.9,0.9,0.9}{
            \makecell[{{p{\linewidth}}}]{
                \texttt{\tiny{[GM$|$GM]}}
                \texttt{Initial states:} \\
\texttt{```} \\
\texttt{Player 1:} \\
\texttt{ 1234567} \\
\texttt{╔═╤═╤═══╗} \\
\texttt{║L│◌│◌◌◌║ 1} \\
\texttt{║C├─┤◌◌◌║ 2} \\
\texttt{╟─┤◌├─┬─╢ 3} \\
\texttt{║◌│◌│◌│◌║ 4} \\
\texttt{╟─┼─┘◌│◌║ 5} \\
\texttt{║◌│◌◌┌┴─╢ 6} \\
\texttt{║◌│◌◌│◌P║ 7} \\
\texttt{╚═╧══╧══╝} \\
\\ 
\\ 
\texttt{Player 2:} \\
\texttt{ 1234567} \\
\texttt{╔═╤═╤═══╗} \\
\texttt{║◌│◌│◌◌◌║ 1} \\
\texttt{║◌├─┤◌L◌║ 2} \\
\texttt{╟─┤P├─┬─╢ 3} \\
\texttt{║◌│◌│◌│◌║ 4} \\
\texttt{╟─┼─┘C│◌║ 5} \\
\texttt{║◌│◌◌┌┴─╢ 6} \\
\texttt{║◌│◌◌│◌◌║ 7} \\
\texttt{╚═╧══╧══╝} \\
\\ 
\texttt{```} \\
            }
        }
    }
    & & \\ \\

    \theutterance \stepcounter{utterance}  
    & & & \multicolumn{2}{p{0.3\linewidth}}{
        \cellcolor[rgb]{0.9,0.9,0.9}{
            \makecell[{{p{\linewidth}}}]{
                \texttt{\tiny{[GM$|$GM]}}
                \texttt{End states:} \\
\texttt{```} \\
\texttt{Player 1:} \\
\texttt{ 1234567} \\
\texttt{╔═╤═╤═══╗} \\
\texttt{║◌│L│◌◌◌║ 1} \\
\texttt{║C├─┤◌◌◌║ 2} \\
\texttt{╟─┤◌├─┬─╢ 3} \\
\texttt{║◌│◌│◌│◌║ 4} \\
\texttt{╟─┼─┘◌│◌║ 5} \\
\texttt{║◌│◌◌┌┴─╢ 6} \\
\texttt{║◌│◌◌│◌P║ 7} \\
\texttt{╚═╧══╧══╝} \\
\\ 
\\ 
\texttt{Player 2:} \\
\texttt{ 1234567} \\
\texttt{╔═╤═╤═══╗} \\
\texttt{║◌│◌│◌◌◌║ 1} \\
\texttt{║L├─┤◌◌◌║ 2} \\
\texttt{╟─┤P├─┬─╢ 3} \\
\texttt{║◌│◌│◌│◌║ 4} \\
\texttt{╟─┼─┘◌│C║ 5} \\
\texttt{║◌│◌◌┌┴─╢ 6} \\
\texttt{║◌│◌◌│◌◌║ 7} \\
\texttt{╚═╧══╧══╝} \\
\\ 
\texttt{```} \\
            }
        }
    }
    & & \\ \\

    \theutterance \stepcounter{utterance}  
    & & & \multicolumn{2}{p{0.3\linewidth}}{
        \cellcolor[rgb]{0.9,0.9,0.9}{
            \makecell[{{p{\linewidth}}}]{
                \texttt{\tiny{[GM$|$GM]}}
                \texttt{* success: True} \\
\texttt{* lose: True} \\
\texttt{* aborted: False} \\
\texttt{{-}{-}{-}{-}{-}{-}{-}} \\
\texttt{* Shifts: 1.00} \\
\texttt{* Max Shifts: 4.00} \\
\texttt{* Min Shifts: 2.00} \\
\texttt{* End Distance Sum: 14.60} \\
\texttt{* Init Distance Sum: 15.76} \\
\texttt{* Expected Distance Sum: 12.57} \\
\texttt{* Penalties: 9.00} \\
\texttt{* Max Penalties: 8.00} \\
\texttt{* Rounds: 12.00} \\
\texttt{* Max Rounds: 12.00} \\
\texttt{* Object Count: 3.00} \\
            }
        }
    }
    & & \\ \\

\end{supertabular}
}

\end{document}
